% !PRINTED 197
\weberpage{197}
%<<CONT
have set) and the upper or lower sign holds according as
$\lambda$ is even or odd. In the latter case this expression is
divisible by $\xi^3 - \eta^3$, and the quotient can be expressed rationally
and in integers through $\xi^3 + \eta^3$ and $\xi\eta$ (because a
symmetric function of two quantities can be represented in this way
through the sum and the product). If we therefore set $\tau = 0$
or $= 1$ according as $\lambda$ is even or odd, there results
for $C'$ an expression of the form
\[
C'(\xi, \eta) = (\xi^3 - \eta^3)^\tau \; \Sigma \, M (\xi\eta)^\alpha (\xi^3 + \eta^3)^\beta,
\]
in which the $M$ are integer coefficients and $\alpha$, $\beta$ run through a series
% [sic] The print says positive integers; as exponents alpha and beta may also be 0, so nonnegative integers are meant.
of positive integers which satisfy the condition
\weq{3}{3\tau + 3\beta + 2\alpha = \nu}
From (1) there results, then, if we multiply by a sufficiently high power of 3 and express $\xi\eta$, $\xi^3+\eta^3$, $\xi^3-\eta^3$
by $H$, $f$, $Q$ according to \S.\,62, (10), and $r$ by $D^{-\frac{1}{6}}$,
\weq{4}{3^\sigma C = Q^\tau \, \Sigma \, M D^\gamma H^\alpha f^\beta,}
in which
\weq{5}{6\gamma = \lambda - 3\tau - 2\alpha,}
and in which the $M$ are likewise integer factors.

To each value of $\gamma$ there belongs, by (5), a definite
value of $\alpha$, and by (3) a definite value of $\beta$; likewise
the two remaining numbers are in each case determined by a value
of $\alpha$ or of $\beta$, so that in (4) every exponent of
$D$, $H$ or $f$ occurs only once.

That $\gamma$ is an integer follows easily from (2), (3), (5);
for since by the definition $\lambda - 3\tau$ is an even number,
$3\gamma$ is first of all an integer, and since further, by (3),
\[
\lambda - 2\alpha = \frac{3}{2}(\mu - \nu + 2\beta + 2\tau),
\]
$6\gamma$, and hence $3\gamma$, is divisible by 3. From this it follows that
% [sic] The print names sigma here. The immediately preceding divisibility argument actually proves that gamma is integral; sigma had already been introduced as a sufficiently high exponent of 3.
$\sigma$ too is an integer.

That $\gamma$ must, however, also be positive, we see as follows:

Suppose that on the right-hand side of (4) negative powers
of $D$ also occur; then we multiply this equation
on both sides by so high a power of $D$ that on the right no
negative powers of $D$ appear any longer, but that the
right-hand side also does not acquire the factor $D$. If in the equation
so obtained we set
%CONT>>
